% Integration notes: theorem statements and proofs, without a document preamble.
% Suggested bibliography: Agol--Hass--Thurston 2006; retained Lackenby July 2026
% manuscript, plus source repository synthesis/parallelity.tex.

\section{Recovering the attaching order of compressed parallelity bundles}
\label{sec:port-order}

The existing source parallelity operator retains complete polygon charts, an
inherited core with at most $30T$ handles, and at most $q\leq66T$ labelled
core ports.  Its component profiles contain Euler characteristic, boundary
count, fibre twisting and a set of port labels.  These are not yet an explicit
small attaching diagram.  In particular, a set of ports does not specify their
cyclic order.  We give an exact extraction of this order and then a relative
replacement theorem for disc products.  This operation fills a specific
geometric gap between a compressed cut and a small patterned handle structure;
it does not assert that the resulting handle structure is admissible.

\subsection{The marked boundary model}

Let $B$ be a finite, possibly disconnected, compact surface supplied as
polygons with pairwise glued sides.  We assume that the quotient really is a
surface, including the interval-or-circle condition on every vertex link.
Each unglued side lies on $\partial B$.  Select $q$ unglued side occurrences,
with distinct labels $0,\ldots,q-1$, called \emph{ports}.  Every port carries
an ordered endpoint chart $(j,0),(j,1)$.  Endpoints of different ports may
coincide after quotienting; even the two endpoints of one boundary edge cell
may coincide.  The statement concerns labelled edge parametrisations, and
does not assume that the closure of each edge cell has distinct endpoints.
All other boundary edges are called ordinary.  In the source application,
ordinary edges lie on the original manifold boundary.

The polygon families and their gluings are represented by affine interval
pairings.  Let $N$ be the number of corner occurrences before quotienting,
and let $R_V$ be the interval equivalence relation that identifies corners
across glued polygon sides.  Thus the $R_V$-orbits are the vertices of $B$.
Let $R_\circ$ be generated by $R_V$ and the endpoint pairings of every
ordinary boundary edge, \emph{omitting the port edges}.  In a compressed
family the omission of a singleton port can be achieved by splitting that
family at its rank, so $q$ omissions create only $O(q)$ additional records.
The source packet already lists each port as a singleton.

\begin{theorem}[Punctured boundary matching]
\label{thm:punctured-boundary-matching}
Give corner occurrence $(j,e)$ the weight $2^{2j+e}$ in the first coordinate.
Give exactly one endpoint occurrence of each ordinary boundary edge weight
$1$ in a second coordinate.  Weights at coincident occurrences add.  Then the
weighted orbit histogram for $R_\circ$ determines all boundary circles
containing ports, their cyclic port order, the entry endpoint of each port,
and the number of ordinary edges in each gap between consecutive ports.
No polygon or boundary-edge family needs to be expanded.

More precisely, every orbit with nonzero first coordinate has weight
\[
  \bigl(2^a+2^b,\,\ell\bigr),\qquad a\ne b,\quad\ell\geq0,
\]
and multiplicity one.  These orbits give a perfect matching $\tau$ on the
$2q$ endpoint tokens.  Let $\iota(2j+e)=2j+1-e$ pair the ends of each port.
Alternating $\iota$ and $\tau$ gives exactly the boundary circles containing
ports.  The matching edge of $\tau$ carries the intervening length $\ell$.
\end{theorem}

\begin{proof}
First quotient by $R_V$.  The free edges form the ordinary boundary
1-manifold of $B$, a disjoint union of circles.  Remove the interiors of its
$q$ port edges.  Each boundary circle containing a port becomes a union of
closed paths, possibly reduced to a single vertex.  Every such path has
exactly two incident endpoint \emph{occurrences} from the removed edges.
This remains true when consecutive ports share a vertex or a one-edge
boundary circle loses its only edge.  A circle without ports remains a
circle and carries no endpoint mark.  Interior vertices of $B$ also carry
no endpoint mark.

The connected components just described are precisely the marked
$R_\circ$-orbits.  Each endpoint bit was placed once, so a marked orbit has
exactly the two stated bits and multiplicity one.  Its second coordinate
counts ordinary boundary edges because every such edge contributes exactly
one unit to its component.  Hence this coordinate is zero exactly for a
zero-length gap; it is positive for an actual interval of ordinary boundary.
Restoring each deleted port joins its two endpoints by $\iota$, and restores
the original circles.  Walking these circles reads their port order and the
direction of each endpoint chart.  This is precisely the alternating walk
in the union of the two perfect matchings.
\end{proof}

There are two directed alternating cycles on endpoint tokens for each
unoriented boundary circle.  The implementation visits each port once, rather
than counting both cycles.  Its canonical choice starts at the least port
label on a circle, enters that port at endpoint $0$, and follows the resulting
direction.  The endpoint labels fix this direction even when the underlying
surface has no selected orientation.  With $q=0$ the output is the empty
list.  Unmarked boundary circles continue to be counted by the source
surface-classification packet, and are not spuriously returned as marked
circles.

For verification it is insufficient to compare only the sum of all endpoint
integers: duplicated bits can carry in binary arithmetic.  The checker
requires two set bits in each nonzero mask, multiplicity one, pairwise
disjoint endpoint-token sets, and their union equal to all $2q$ tokens.
The underlying source-bound construction also authenticates that each mark
was placed at exactly its specified source corner occurrence.

\subsection{Relative fibre charts}

Let $E\to B$ be the source $I$-bundle.  Each polygon has a local product
chart, and every side gluing specifies whether it reverses the fibre.
These signs define the fibre-orientation double cover $\widetilde B\to B$.
The existing source packet checks that they arise from the actual source
bundle.  In particular, they are compatible around vertices; arbitrary
signs on a polygon adjacency graph would not suffice to define a bundle.
In these source product charts the fibre transition maps are exactly
$t\mapsto t$ and $t\mapsto1-t$.  Thus the recorded gauge bits determine
the parametrised fibre charts themselves, not merely an abstract interval
bundle isomorphism class with arbitrary fibre reparametrisations.

\begin{theorem}[Marked fibre gauges]
\label{thm:marked-fibre-gauges}
Place a distinct bit $2^{2j+e}$ at sheet $e$ of a polygon occurrence incident
to port $j$ in $\widetilde B$.  One weighted orbit query determines, for
each port-bearing component of $B$, whether the fibre-orientation cover is
connected.  When $E$ is a product over that component, it also determines
the relative fibre reversal at every port, after the gauge at any one
anchor port has been fixed.
\end{theorem}

\begin{proof}
Over a connected component of $B$, the orientation cover has two connected
components precisely when the bundle has a global fibre orientation.  In
that case each cover component contains exactly one of the two marks of
every port.  Select the component containing sheet $0$ of the anchor port.
If it contains sheet $e_j$ of port $j$, the local chart at $j$ differs from
the anchor trivialisation by the reversal bit $e_j$.  The other component
contains the complementary marks.  If the cover is connected, it contains
both marks of every port and no such global gauge exists.  Unique port
labels identify these components without any multiplicity ambiguity.
\end{proof}

This query is independent of the cyclic-order query.  Reversing a base edge
parametrisation and reversing the interval fibre are different operations.
In particular, a singleton rank map does not encode fibre reversal.  The
implementation retains both the signed entry endpoint and the fibre bit.

\subsection{A relative disc-product replacement theorem}

Suppose a port-bearing component of $B$ is a disc $D$.  Write $A$ for the
union of its port edge cells, including their endpoints.  Let $E\to D$ be
the product bundle and let $K$ be the complementary core.  Its attachment
to $K$ is the specified source map from $E|_A$; endpoints shared by port
cells retain their identifications and their source charts.
An ordinary gap is unlabelled for this replacement, except for its status
as original boundary and the retained numerical count of its source edges.
The replacement does not preserve the discarded interior subdivision of
such a gap as additional marked data.

\begin{theorem}[Relative disc-product model]
\label{thm:relative-disc-product-model}
Let a disc component have $q_D\geq1$ port cells.  Its signed cyclic word,
zero-versus-positive gap data and relative fibre gauges determine a product
model with at most $9q_D$ tetrahedra, together with all port attachment charts.
Replacing $E$ by this model in $K\cup E$ gives a PL homeomorphic manifold
by a homeomorphism restricting to the identity on $K$.

The same homeomorphism preserves the first-cut boundary pattern when the
pre-cut pattern was empty: within $E$ this pattern is the two horizontal
rims over the ordinary part of $\partial D$.  It makes no assertion about
an arbitrary unrecorded inherited pattern.
\end{theorem}

\begin{proof}
Let $r$ be the number of positive ordinary gaps.  Retain one side for every
port and one side for every positive gap, inserting no side at a zero gap.
Thus the new boundary cell structure has
\[
 m=q_D+r\leq2q_D
\]
sides.  It has exactly the original cyclic order of labelled ports and
exactly the original endpoint coincidences between consecutive ports.
When $m<3$, subdivide one side so that there are $n=3$ sides; otherwise put
$n=m$.  Each subdivision carries its rational parameter interval in the
original side chart.  Hence $n=\max\{3,m\}\leq3q_D$.

Fill the boundary polygon by a fan from one interior vertex, giving $n$
triangles.  The product of each triangle with $I$ has a staircase
triangulation into three tetrahedra.  Use one global total order on the
base vertices; the induced diagonals then agree on shared rectangular
faces.  This gives $3n\leq9q_D$ tetrahedra.

The signed cyclic word supplies a PL homeomorphism from the model boundary
circle to $\partial D$ that agrees with every port parametrisation.  On
positive ordinary gaps choose any monotone PL extension between the given
endpoint images.  At zero gaps the common endpoint is retained.  The map
extends to a PL homeomorphism of discs.  This is the elementary disc
extension theorem; equivalently, triangulate both discs compatibly with the
chosen boundary map and extend over their interiors.

Since a disc is contractible, its interval bundle is a product.  Choose a
global fibre trivialisation whose local gauge is the prescribed sheet $0$
at the anchor port.  The marked cover theorem forces its restriction at
every other port to be exactly the stored fibre bit.  Taking the product
with the extended disc homeomorphism therefore gives a bundle homeomorphism
that agrees with every specified source rectangle chart.  On a model side
with parameter $u$ and vertical parameter $t$, that agreement uses
\[
\begin{aligned}
 u&\longmapsto u\ \hbox{or}\ 1-u
 &&\text{according to the entry endpoint},\\
 t&\longmapsto t\ \hbox{or}\ 1-t
 &&\text{according to the fibre gauge}.
\end{aligned}
\]
Restrict these maps to the stored rational subintervals when a side was
subdivided.  Glue this homeomorphism to the identity on $K$.  The maps agree
on every attaching chart, including coincident endpoint fibres, so the
quotient map is a PL homeomorphism of the unions.

For an initially empty pattern, the new pattern inside the product is
exactly $\overline{\partial D\setminus A}\times\partial I$, with the
appropriate endpoint convention at the port interfaces.  The constructed
map preserves the ordinary boundary intervals and maps the two horizontal
sheets according to the recorded gauges; hence it preserves this pattern.
Any further marked data requires its own preserved charts.  If $r$ extra
boundary subdivision points suffice to carry such data in product charts,
the same construction uses $O(q_D+r)$ tetrahedra; one cannot retain a bound
independent of $r$ merely by omitting those marks from the input.
\end{proof}

\begin{corollary}[Linear total disc-product template size]
For a source triangulation with $T$ tetrahedra, all port-bearing disc
components of its source parallelity bundle can simultaneously be replaced
using at most
\[
 9q\leq594T
\]
tetrahedra in total, with $O(T)$ labelled port records.  The core remains
the same source-bound core with at most $30T$ inherited handles.
\end{corollary}

\begin{proof}
Every port belongs to exactly one component, so the sum of $q_D$ over these
disc components is at most $q$.  The preceding relative homeomorphisms
restrict to the identity on $K$ and therefore apply simultaneously.  The
source bounds $q\leq66T$ and the $30T$ core-handle bound give the estimates.
\end{proof}

This corollary is a bound on the added disc templates.  It is not itself a
tetrahedron count for the entire cut: non-disc bundle components may remain,
and the core must still be triangulated compatibly with its source attaching
charts.  Port-free disc bundles are separate components from the core and
can be retained as counted product types.  Explicitly listing an enormous
number of such balls would defeat compression.  Most importantly, the
operation replaces a disc product; it does not delete it.  A ball attached
along several rectangles need not be removable without changing topology.

\subsection{Why the extra boundary data is necessary}

\begin{proposition}[Port profiles do not determine marked base attachments]
The Euler characteristic, number of boundary components, orientability,
set of port labels, and product/twisting type do not determine the labelled
base of a disc product.  Consequently they do not determine it in the
category of bundle homeomorphisms preserving the labelled port charts.
\end{proposition}

\begin{proof}
Take a disc with four pairwise disjoint labelled port arcs.  In one model
their cyclic order is $(1,2,3,4)$; in another it is $(1,3,2,4)$.  Give the
ordinary gaps equal subdivisions and choose trivial fibre charts.  All
listed numerical data and port sets agree.  A homeomorphism of discs
restricts to a homeomorphism of their boundary circles, which preserves or
reverses cyclic order.  The two labelled cyclic words are not related by a
rotation or reversal.  Hence there is no homeomorphism preserving all port
labels.  A bundle homeomorphism induces a base homeomorphism and therefore
cannot identify these two marked products.
\end{proof}

The category matters here: without the fibre structure or further boundary
marking, an arbitrary homeomorphism of a 3-ball can permute disjoint disc
patches on its boundary.  The cyclic-order argument alone does not rule out
such a homeomorphism.  If the first-cut rim pattern is also retained and all
port fibre gauges are fixed to zero, that pattern supplies an additional
obstruction: its arcs join exactly the adjacent port endpoints.  The two
four-port examples have different labelled adjacency relations and cannot
be identified while preserving both those rim arcs and the port charts.

Even a cyclic word needs zero-gap information.  Two adjacent ports have a
common endpoint fibre, whereas two separated ports do not.  Their signed
two-port words can agree.  The total number of ordinary edges can also be
made equal by subdividing the unique ordinary gap in the adjacent case.
The second coordinate in Theorem~\ref{thm:punctured-boundary-matching}
retains the gap-by-gap distinction and prevents this loss of attaching data.

\subsection{Encoded complexity and an unconditional orbit-count reduction}

The practical extractor uses one two-coordinate weighted query for port
paths and one scalar weighted query for fibre marks, in addition to the
four existing source queries.  If the polygon packet has $s$ interval
records, the two new query inputs have $O(s+q)$ records.  Endpoint masks
have $2q$ bits; this matters for bit complexity.  The marking input can use
$O(q^2)$ bits even though it has only $O(q)$ records.  With
$L=\lceil\log_2(N+2)\rceil$ and the analogous logarithm for the polygon
universe, a coarse input bound is
\[
 O\bigl((s+q)L+q^2\bigr).
\]
The output contains $O(q)$ port and template records.  Gap counts use
$O(L)$ bits each.  The supplied implementation's weighted replay has its
existing dependence on the realized orbit trace and weight-run partitions.
We do not infer a uniform polynomial bound for those partitions from a
small query description.

There is nevertheless a separate deterministic polynomial construction of
the \emph{topological} port data, using only unweighted interval-orbit
counting.  This gives a useful distinction between the mathematical
availability of a bounded extraction and the specific faster weighted
implementation.

\begin{lemma}[Orbit-membership by a singleton merge]
Let $c(R)$ be the number of orbits of an interval-pairing relation $R$ on a
finite universe, and let $[x\sim y]$ denote one added singleton pairing.
Then
\[
 c\bigl(R\cup[x\sim y]\bigr)=
 \begin{cases}
 c(R),&x,y\text{ already belong to the same orbit},\\
 c(R)-1,&\text{otherwise}.
 \end{cases}
\]
\end{lemma}

\begin{proof}
The added relation identifies exactly the two existing equivalence classes
containing $x$ and $y$.  They are either the same class or distinct classes.
\end{proof}

\begin{lemma}[Orbit membership in a set by a singleton star]
\label{lem:orbit-star-membership}
Let $R$ be an interval-pairing relation and let $U$ be a nonempty finite
list of marked positions.  Choose $u_0\in U$, and let $R[U]$ be obtained
from $R$ by adding singleton pairings $u_0\sim u$ for all $u\in U$.
Then, for any position $x$,
\[
 c\bigl(R[U]\cup[x\sim u_0]\bigr)=c(R[U])
 \quad\Longleftrightarrow\quad
 [x]_R\cap U\ne\varnothing.
\]
The construction uses at most $|U|-1$ nontrivial added pairings, and remains
valid when distinct endpoint tokens occupy the same position.
\end{lemma}

\begin{proof}
The singleton star merges exactly the original $R$-classes that contain a
position of $U$, and changes no other class.  Hence $x$ is in the resulting
class of $u_0$ exactly when its original class meets $U$.  Apply the
singleton-merge lemma to $R[U]$.  Repeated positions merely give redundant
or identity pairings and do not change this reasoning.
\end{proof}

\begin{theorem}[Unweighted polynomial extraction]
For a valid compressed polygon surface with $q$ labelled boundary ports,
their boundary matching, signed cyclic words and zero-gap flags can be
computed with $O(q\log(q+1))$ ordinary interval-orbit counting calls.  Product
components' relative fibre gauges require at most $O(q^2)$ further calls,
or $O(q)$ when the port partition into base components is supplied.  Each
call has $O(s+q)$ interval records and the same binary universe scale.
Consequently these data have a deterministic polynomial-time extraction in
$s,q$ and the binary universe length, by the Agol--Hass--Thurston orbit
counting theorem.  With the supplied source component partition, all these
topological chart data use $O(q\log(q+1))$ unweighted counting calls in total.
\end{theorem}

\begin{proof}
Choose the least unmatched endpoint token $x$.  By
Theorem~\ref{thm:punctured-boundary-matching}, exactly one other unmatched
endpoint token is in its $R_\circ$-class.  Bisect the remaining candidate
list.  For a nonempty half $U$, apply
Lemma~\ref{lem:orbit-star-membership} to $R_\circ$, using two ordinary orbit
counts.  Since $x$ is not a token of $U$, the answer says precisely whether
its unique partner is in that half.  Continue with the indicated half until
one candidate remains, remove the matched pair, and repeat.  Each test starts
with the original relation $R_\circ$ and its own star; stars from earlier
tests are not retained.  There are $O(q\log(q+1))$ tests altogether and at
most $O(q)$ extra singleton pairings in any one test.  Coincident marked
positions are handled by the same lemma.  This recovers $\tau$.
For each matched pair, repeat the membership test using $R_V$ instead of
$R_\circ$.  Its gap has zero ordinary edges if and only if its endpoints
are the same quotient vertex.  Otherwise the gap is a nonempty path.
These $q$ further tests use the simpler singleton-merge lemma and a shared
baseline count for $R_V$.  Reading the alternating matching cycles gives
the signed words.

If required, partition port polygon occurrences into base components using
the same membership test on the polygon-gluing relation, in $O(q^2)$ calls.
For each component choose an anchor.  Its two sheet marks lie in the same
fibre-cover orbit exactly in the twisted case.  In the product case,
compare each port's sheet-$0$ occurrence with the anchor's sheet-$0$
occurrence.  Equality gives gauge $0$; inequality gives gauge $1$, since
exactly one of that port's sheets lies in the anchor component.  These are
$O(q)$ more tests.

All added pairings are singletons, with at most $O(q)$ of them in any
instance.  Apply the ordinary polynomial interval-orbit counting algorithm
to each instance.  The number of calls and the encoded size of each instance
are polynomial in the stated parameters.  Numerical gap lengths are not
obtained by this argument; their zero-versus-positive flags are the
information required by the relative topological replacement theorem.
\end{proof}

The unweighted construction is a proved reduction, not the code path
benchmarked here.  The implementation uses weighted extraction for fewer
queries and records exact gap lengths.  This distinction prevents a
trace-sensitive implementation measurement from being silently upgraded
into a hierarchy-wide complexity claim.

\subsection{Certificates, experiments and remaining obligations}

The code files are \texttt{parallelity\_port\_words.py} and
\texttt{parallelity\_port\_words\_verify.py}.  The certificate includes the
entire existing source parallelity certificate, two new independently
replayable weighted-orbit proofs, signed boundary words, fibre gauges, and
disc-product templates.  The verifier first authenticates the source
geometry using the existing independent raw-cell reconstruction.  It then
independently constructs the two queries, replays their proofs, decodes the
boundary and fibre information by separate code, and reconstructs the
template charts and triangulations.  It does not call a port-word producer
helper.  All six discovery queries share the cycle budget.  The verification
cap counts all six supplied trace lengths.  An exhausted discovery returns
no geometry, word, template or certificate; an exhausted replay is false.

The finite oracle in \texttt{parallelity\_port\_research.py} explicitly
glues all polygon corners, materializes boundary edges and follows their
two incident darts.  A separate breadth-first propagation on the polygon
adjacency graph determines fibre gauges.  Thus it checks geometric order,
adjacent ports and fibre data independently of weighted histograms.

The research audit also assembles the emitted tetrahedra with Regina and
checks validity, connectedness, orientability, one boundary component and
the exact ball-recognition predicate.  This is a check of the actual emitted
complexes, beyond an Euler-characteristic or homology comparison.  It is not
a dependency of the production extractor or certificate replay; their
correctness is supplied by the explicit disc-product construction above.
Every port's rational subdivision intervals are separately checked to form
an exact partition of $[0,1]$.

The tests include zero ports, one port, adjacent ports, an actual triangular
disc quotient whose only boundary edge has coincident endpoints, source
self-gluings, relabellings, disconnected bases, one-sided surfaces,
independently replayed hostile certificates, exact binary transport of
$2^{16384}$ source multiplicity, shared caps and cancellation.  The source
two-tetrahedron example with a single selected two-handle has eight adjacent
ports.  Its replacement has eight fan triangles and twenty-four tetrahedra;
there are no original-boundary gaps to insert.

For integration into the hierarchy, the main remaining obligations are:
carry an arbitrary existing pattern through these charts; supply any
pattern marks lying in product horizontal faces; construct and verify the
admissible patterned handle structure; control non-disc bundle components;
and prove the reductions and search bounds across successive hierarchy
stages.  The present homeomorphism theorem preserves a supplied cut but does
not decide whether a decomposition square is trivial or whether a boundary
pattern is essential.  It also does not assert the new small triangulation
has bounded handle type suitable for every later normality condition.

% Editorial source note, separate from the proved extraction:
% The retained manuscript's admissibility clause (2) says a disc in
% partial H^0, without explicitly saying disjoint from F.  Its later claim
% about a non-disc component of partial H^0 cut(F union P) appears to use
% a stronger obstacle-relative disc condition.  Do not silently change the
% definition.  The results above do not require resolving this ambiguity.
